%----------------------------------------------------------------------------------------------------
% START OF CHAPTER
%----------------------------------------------------------------------------------------------------
\chapter{March 7}
\paragraph{Topics} 
\textit{Modal verbs, shopping, recap of verbs}


\section{Grammar}
\subsection{Modal verbs}
  A \textbf{modal verb} expresses sentiments like request, capability, and obligation. 
	Used with the infinitive. Italian modal verbs: \itb{volere} (to want), \itb{potere} 
	(to be able), \itb{dovere} (to have to). E.g. `want to eat' would be a conjugated form of 
	\iti{volere} `want' with the infinitive \iti{mangiare} `to eat.' See Table \ref{tab:tVerbsModal}.  

\begin{table}[ht]
\centering
    \begin{tabular}{ r l  | r l  | r l  }
         \itb{potere} & \iti{to be able} & \itb{volere} & \iti{to want} & \itb{dovere} & \iti{to have to} \\
         \hline
         \itb{io}   & \iti{posso}      & \itb{io}      & \iti{voglio}   & \itb{io}   & \iti{devo} \\
         \itb{tu}   & \iti{puoi}     & \itb{tu}      & \iti{vuoi}     & \itb{tu}   & \iti{devi} \\
         \itb{lui/lei} & \iti{può}   & \itb{lui/lei} & \iti{vuole}    & \itb{lui/lei} & \iti{deve} \\
         \itb{noi}  & \iti{possiamo} & \itb{noi}     & \iti{vogliamo} & \itb{noi}  & \iti{dobbiamo} \\
         \itb{voi}  & \iti{potete}   & \itb{voi}     & \iti{volete}   & \itb{voi}  & \iti{dovete} \\
         \itb{loro} & \iti{possono}  & \itb{loro}    & \iti{vogliono} & \itb{loro} & \iti{devono}

    \end{tabular}
    \caption{Modal verbs, present tense}
    \label{tab:tVerbsModal}
\end{table}
\paragraph{Examples}
\begin{itemize}
    \item I want a coffee. \\
    \iti{Io voglio un caffè.}

    \item I want to buy a coffee. \\
    \iti{Io voglio comprare un caffè.}

    \item I can pay. \\
    \iti{Io posso pagare.}

    \item I have to pay. \\
    \iti{Io devo pagare.}
\end{itemize}

\subsection{Conditional usage}
Just as in English, modal verbs are commonly expressed in a conditional form 
(``could, would, should''). E.g., it is more polite to say, ``I would like'' than ``I want''.

\paragraph{Examples}
\begin{itemize}
    \item I would like a coffee. \\
    \iti{Io vorrei un caffè.}

    \item I would like to buy a coffee. \\
    \iti{Io vorrei comprare un caffè.}

    \item Could I pay? \\
    \iti{Io potrei pagare?}

    \item Should I pay? \\
    \iti{Io dovrei pagare?}
\end{itemize}

\section{Vocabulary}
\subsection{Words}
\begin{table}[ht] %[ht] % place table roughly [h]ere or else at the [t]op of page
\centering
    \begin{tabular}{ r l  | r l  | r l  }
         store          & \iti{negozio}    & food store    & \iti{supermercato} & brand         & \iti{marca} \\
         restaurant     & \iti{ristorante} & pharmacy      & \iti{farmacia}     & post office   & \iti{ufficio postale} \\
         closed         & \iti{chiuso}     & open          & \iti{aperto}       & hotel         & \iti{albergo} \\
         sale           & \iti{saldi}      & discount      & \iti{sconto}       & receipt       & \iti{scontrino}
	    \tablefootnote{It is expected practice to ask for a receipt for a transaction, 
	    and in fact you may be asked to produce one after leaving a business by the authorities.} \\
         small          & \iti{piccolo}    & medium        & \iti{medio}        & large         & \iti{grande} \\
         suit (clothes) & \iti{abito}      & dress         & \iti{vestito}      & pair of shoes & \iti{paio di scarpe} \\
         on sale        & \iti{in offerta} & special offer & \iti{offerta speciale} \\
         size (clothes) & \iti{taglia}     & size (shoes)  & \iti{misura}    
    \end{tabular}
    \caption{Shopping terms}
    \label{tab:tVocabShopping}
\end{table}

\subsection{Phrases}

\begin{itemize}
    \item Io parlo italiano ma io sono un principiante. \\
      \iti{I speak Italian, but I'm a beginner.}

    \item Puoi ripetere? \\
      \iti{Can you repeat?}

    \item Per favore, parla lento. \\
      \iti{Please, speak slowly.}

    \item Io parlo un po' italiano. \\
      \iti{I speak a little Italian.}
      \footnote{Contrast \iti{po'}, expressing a limit, to \iti{poco}, representing quantity (`few').}
\end{itemize}

\section{Conversation: \iti{Al negozio}}

\begin{itemize}
    \item [Luigi] Good morning. Can I help you? \\
        \iti{Buongiorno. Posso aiutare?}
    \item [Francesca] Hello. Yes. I want to buy a pair of shoes, please. \\
        \iti{Salve. Sì. Io vorrei comprare un paio di scarpe, per favore.}
    \item [Luigi] What size? \\
        \iti{Quale misura?}
    \item [Francesca] Forty, thanks. I need to also buy an elegant dress for this evening. \\
        \iti{40, grazie. Io devo comprare anche un vestito elegante per questa sera.}
    \item [Luigi] What size? \\
        \iti{Quale taglia?}
    \item [Francesca] Smaller. Medium please. How much does it cost?    \\
        \iti{Più piccolo. Media per favore. Quanto costa?}        
\end{itemize}

\section{Homework}
\begin{enumerate}
    \item I am very happy because I am eating Claudio's pizza. \\
      \iti{Io sono molto felice perché sto mangiando la pizza di Claudio.}
    
    \item Claudio is not happy because I am eating Claudio's pizza. \\
      \iti{Claudio non è felice, perché io sto mangiando la sua pizza.}
    
    \item Your glass of wine is on the table. \\
      \iti{Il tuo bicchiere di vino è sul tavolo.} 
    
    \item His bottle of wine is very expensive. \\
      \iti{La sua bottiglia di vino è molto cara.}    
    
    \item We are studying Italian with Claudio's book. \\
      \iti{Noi stiamo studiando italiano con il libro di Claudio.}
    
    \item My wallet is in the car. \\
      \iti{Il mio portafoglio è nella macchina.}
    
    \item This bread is very good because it is fresh. \\
      \iti{Questo pane è molto buono perché è fresco.}
    
    \item I am speaking with Claudio. \\
      \iti{Io sto parlando con Claudio.}
    
    \item I am sick. (Using the `to stay' verb) \\
      \iti{Io sto male.}
    
    \item I am sick. (Using the `to be' verb) \\
      \iti{Io sono malato.}          
\end{enumerate}
%----------------------------------------------------------------------------------------------------
% END OF CHAPTER
%----------------------------------------------------------------------------------------------------

